\chapter{Il dataset}
\label{cap:dati}

Questo capitolo descrive i \textbf{dati} su cui si fonda l'intero lavoro. Si parte dai
dati grezzi --- le simulazioni idrodinamiche del campo di concentrazione e i dati
meteo-oceanografici di vento e corrente (Sezione~\ref{sec:dati}) --- per poi illustrare
come questi vengono caricati, sincronizzati e gestiti via software
(Sezione~\ref{sec:dataloader}). Su questi dati viene costruito l'ambiente di simulazione
vero e proprio, oggetto del Capitolo~\ref{cap:ambiente} e condiviso dai due approcci
confrontati nella tesi: ci\`o che cambia tra i due \`e unicamente la politica che governa
l'agente, non i dati n\'e l'ambiente.

\section{Dati utilizzati}
\label{sec:dati}

\subsection{Modello idrodinamico MIKE21}
\label{sec:mike}

Il campo di concentrazione del contaminante \`e ottenuto da simulazioni idrodinamiche
realizzate con \textbf{MIKE21 Flow Model FM} (\emph{Danish Hydraulic Institute},
DHI)~\citep{mike21}, un codice ai volumi finiti basato sulla formulazione bidimensionale
delle equazioni di Navier--Stokes per flussi incomprimibili che, sotto le ipotesi di
Boussinesq e di distribuzione idrostatica della pressione, si riduce alle equazioni delle
acque basse (\emph{Shallow Water Equations}). Il codice include la simulazione dei
processi di advezione e dispersione/diffusione di un inquinante non reattivo, trattato
con approccio euleriano: il trasporto \`e descritto a partire da sorgenti di massa
caratterizzate da una portata e da una concentrazione assegnate.

La discretizzazione spaziale del dominio \`e realizzata mediante una \textbf{griglia a
elementi triangolari di dimensione variabile}, che consente un'elevata risoluzione nelle
aree di maggiore interesse mantenendo una discretizzazione pi\`u grossolana altrove, al
fine di contenere i tempi di calcolo. Il dominio di calcolo comprende l'area circostante
il porto di Cecina (Figura~\ref{fig:dominio}), con un'estensione di circa
$16 \times 14\,\text{km}$ e una dimensione
caratteristica della griglia variabile fra circa $10\,\text{m}$ in prossimit\`a del porto
(zona di interesse) e circa $150\,\text{m}$ nelle aree pi\`u lontane.

\begin{figure}[H]
  \centering
  \includegraphics[width=\linewidth]{dominio_batimetria.png}
  \caption{Dominio di calcolo del modello idrodinamico. A sinistra: batimetria e mesh a
           elementi triangolari di dimensione variabile (pi\`u fitta nella zona di
           interesse). A destra: estensione del dominio ($\sim$$16\times14\,\text{km}$),
           area di interesse (riquadro) da cui \`e estratto il dataset e nodi CMEMS lungo i
           bordi usati per le condizioni al contorno. Baia di Cecina (costa toscana).}
  \label{fig:dominio}
\end{figure}

\subsection{Forzanti e condizioni al contorno}
\label{sec:forzanti}

Il modello locale realizza un \emph{downscaling idrodinamico} del bacino del
Mediterraneo, incrementando la risoluzione spaziale da circa $4{,}5\,\text{km}$ del
modello a grande scala fino a circa $10\,\text{m}$ nell'area di interesse. Le condizioni
al contorno di livello della superficie libera e di velocit\`a della corrente sono
derivate dal prodotto CMEMS~\citep{cmems}
\texttt{MEDSEA\_\allowbreak ANALYSIS\_\allowbreak FORECAST\_\allowbreak PHY\_\allowbreak 006\_013},
che fornisce dati orari delle componenti zonale e meridionale della velocit\`a di
corrente e del livello marino. Per ciascun nodo di calcolo ai bordi del dominio sono
imposte condizioni al contorno di tipo \emph{Flather}~\citep{oddo_pinardi}, mediante
l'imposizione dei valori orari delle componenti $(u, v)$ e dell'elevazione della
superficie libera, interpolati linearmente dai nodi CMEMS disponibili lungo i bordi.

La forzante di vento \`e imposta come variabile nel tempo e costante sul dominio, ed \`e
estratta dal dataset di \emph{reanalysis} \textbf{ERA5}~\citep{era5} dell'ECMWF
(componenti del vento a $10\,\text{m}$, $u_{10}, v_{10}$, a risoluzione oraria). In questa
fase non vengono considerati alcuni contributi al trasporto assunti come secondari: le
correnti indotte dal moto ondoso (deriva litoranea, correnti di ritorno) e gli apporti di
acqua dolce (scarichi fluviali, runoff superficiale).

\subsection{Scenari idrodinamici e di vento}
\label{sec:scenari}

Le serie temporali CMEMS ai bordi del dominio (2024--2025) sono state classificate in
finestre di 48 ore tramite descrittori sintetici (intensit\`a media della corrente,
persistenza direzionale, variabilit\`a del livello marino) e raggruppate con un algoritmo
di \emph{clustering} k-means in 4 classi idrodinamiche di riferimento, che rappresentano
regimi distinti in termini di intensit\`a della corrente, persistenza direzionale e
variabilit\`a del livello: dalla Classe~1 (regime moderato con elevata variabilit\`a del
livello) alla Classe~4 (regime energetico e altamente persistente). Le simulazioni
adottate in questa fase si riferiscono alla sola \textbf{Classe 2} (regime moderato e
persistente, con correnti di intensit\`a intermedia, direzione stabile e bassa
variabilit\`a del livello, ovvero condizioni relativamente stazionarie e ben organizzate),
rappresentata dalla finestra del 19--20 maggio 2024.

Analogamente, l'analisi statistica delle serie di vento ERA5 (intensit\`a, direzione,
persistenza; Figura~\ref{fig:rosa_venti}) ha portato a definire \textbf{4 classi di vento}
rappresentative, sovrapposte alla condizione idrodinamica di riferimento:
\begin{itemize}
    \item \textbf{V0} --- \emph{condizioni di calma / vento debole}. Vento di bassa
          intensit\`a (circa 1{,}5--2\,m/s), senza una direzione dominante significativa;
          rappresenta condizioni con forcing eolico trascurabile.
    \item \textbf{V1} --- \emph{vento persistente da NO}. Direzione prevalente compresa tra
          310\textdegree{} e 330\textdegree, con velocit\`a
          fino a 8--10\,m/s; rappresenta uno dei regimi dominanti lungo costa.
    \item \textbf{V2} --- \emph{vento persistente da SO}. Direzione prevalente tra
          230\textdegree{} e 250\textdegree, con velocit\`a
          fino a 8--10\,m/s; rappresenta il regime opposto rispetto al precedente, con
          effetti rilevanti sull'inversione del trasporto lungo costa.
    \item \textbf{V3} --- \emph{vento da settore orientale}. Direzione prevalente tra
          110\textdegree{} e 130\textdegree, con velocit\`a
          fino a 6--8\,m/s; rappresenta un regime meno frequente ma significativo, associato
          a condizioni di trasporto trasversale o secondario.
\end{itemize}

\begin{figure}[H]
  \centering
  \includegraphics[width=0.62\linewidth]{rosa_venti.png}
  \caption{Rosa dei venti dell'area (dati ERA5 2024--2025): distribuzione di frequenza per
           direzione e intensit\`a. I settori dominanti da NO e da SO --- ben visibili ---
           motivano le classi di vento V1 e V2; le classi V0--V3 sintetizzano in modo
           compatto questa climatologia eolica.}
  \label{fig:rosa_venti}
\end{figure}

Per le classi V1--V3 il vento \`e modellato con un profilo temporale tipico: 8 ore di
calma iniziale, 6 ore di aumento lineare fino all'intensit\`a massima, 12 ore di intensit\`a
massima e 6 ore di diminuzione lineare. Ogni scenario genera plume con morfologie
marcatamente diverse (Figura~\ref{fig:plume}): da configurazioni compatte e direzionate a
configurazioni diffuse e frammentate.

\begin{figure}[H]
  \centering
  \includegraphics[width=\linewidth]{plume_scenarios.png}
  \caption{Campo di concentrazione del contaminante (plume) per la stessa sorgente
           ($\star$) sotto i quattro scenari di vento, allo stesso istante (chunk Q1/2).
           In V0 (calma) il plume resta compatto attorno alla sorgente; con vento
           persistente (V1 da NO, V2 da SO) viene avvettato lungo costa e si allunga/diffonde;
           in V3 (da E) \`e sospinto verso il largo. Mare in azzurro, terraferma in bianco,
           concentrazione in scala gialla--rossa.}
  \label{fig:plume}
\end{figure}

\subsection{Sorgenti e formato dei dati}
\label{sec:sorgenti}

Nell'area di interesse sono distribuite \textbf{132 sorgenti} equispaziate di circa
$200\,\text{m}$, ciascuna caratterizzata da una portata costante di $50\,\text{l/s}$ e da
una concentrazione di $1000\,\text{g/m}^3$ dell'inquinante target. Ogni sorgente \`e
simulata sotto i 4 scenari di vento, per un totale di $132 \times 4 = 528$ simulazioni.
I campi di concentrazione risultanti sono estratti in formato \textbf{NetCDF} su una
griglia regolare con passo di $10\,\text{m}$ e intervallo temporale di 2 minuti,
nell'area compresa fra le coordinate $x \in [619000, 622000]$ e
$y \in [4794500, 4797000]$\,m (UTM32N), corrispondente a un riquadro di
$3 \times 2{,}5\,\text{km}$ ($300 \times 250$ celle). I file seguono la convenzione di
nome \texttt{CL02\_<vento>\_SRC<num>\_Conc\_10mGrid.nc}; i corrispondenti campi di
corrente $(u, v)$, comuni a tutte le sorgenti di uno stesso scenario, seguono la
convenzione \texttt{CL02\_<vento>\_U\_V\_10mGrid.nc} (i campi di corrente sono comuni a
tutte le sorgenti di uno stesso scenario, in quanto la presenza della sorgente altera il
campo di moto solo localmente). Sorgente e scenario di ciascun file sono ricavati dal nome
tramite una semplice espressione regolare:

\begin{lstlisting}
m = re.match(r"CL02_(V\d)_(SRC\d+)_Conc", filename)
version, source_id = m.group(1), m.group(2)   # es. ("V2", "SRC042")
\end{lstlisting}

\subsection{Suddivisione del dataset e augmentation temporale}
\label{sec:split}

La suddivisione del dataset e l'augmentation temporale descritte di seguito sono
\textbf{comuni a entrambi gli approcci} confrontati nella tesi (Capitoli~\ref{cap:fcm}
e~\ref{cap:rl}), e fanno quindi parte dell'impostazione condivisa dell'ambiente di
simulazione.

Per garantire una valutazione su sorgenti mai osservate durante l'addestramento del
modello RL, le 132 sorgenti sono suddivise in un insieme di \textbf{training}
(SRC001--SRC106, 80\%, 106 sorgenti) e in un insieme di \textbf{inferenza} held-out
(SRC107--SRC132, 20\%, 26 sorgenti), riservato alla sola valutazione. Le stesse 26
sorgenti held-out sono usate per valutare sia il metodo del gradiente sia il modello RL,
cos\`i da garantire un confronto equo tra i due approcci.

Poich\'e il campo evolve nel tempo, ogni simulazione viene inoltre suddivisa in
\textbf{3 chunk temporali} che corrispondono a diversi stadi della dispersione del plume
e fungono da condizioni iniziali distinte: \textbf{Q1/4} (inizio della propagazione,
plume compatto), \textbf{Q1/2} (stadio intermedio, plume parzialmente disperso) e
\textbf{Q3/4} (stadio avanzato, plume pesantemente disperso). Questa
\emph{augmentation temporale} consente di inizializzare l'agente in tre fasi diverse della
dispersione, aumentando la variabilit\`a delle condizioni. La combinazione di scenari di
vento e chunk temporali definisce le \emph{configurazioni} $(v, q)$ su cui sono
articolati i risultati nei capitoli seguenti. In codice, la suddivisione delle sorgenti e
la mappatura chunk\,$\to$\,frame iniziale si riducono a:

\begin{lstlisting}
sources = sorted(discover_sources())               # SRC001 ... SRC132
train = [s for s in sources if int(s[3:]) <= 106]  # 80%
test  = [s for s in sources if int(s[3:]) >  106]  # 20% held-out

nt = field.n_timesteps                             # numero di frame
spawn_frame = {0: nt//4, 1: nt//2, 2: (nt*3)//4}[chunk_id]   # Q1/4, Q1/2, Q3/4
\end{lstlisting}

\section{Caricamento e gestione dei dati}
\label{sec:dataloader}

Il caricamento dei dati grezzi e la loro trasformazione in strutture interrogabili
dall'ambiente sono gestiti dal modulo software \texttt{data\_loader}, organizzato attorno
a tre componenti principali.

\paragraph{Campo di concentrazione.}
Il componente \texttt{NetCDFLoader} legge un file NetCDF e ne costruisce un oggetto
\texttt{ConcentrationField}, che incapsula i dati di concentrazione $[t, y, x]$ e
fornisce query spazio-temporali continue. Data una posizione $(x, y)$ e un istante $t$,
il campo restituisce la concentrazione tramite interpolazione (bilineare nello spazio,
lineare tra i due frame temporali adiacenti), cos\`i che l'agente non sia vincolato ai
soli centri-cella della griglia. All'inizializzazione, dalla maschera di terra viene
inoltre precalcolata --- tramite \emph{Euclidean Distance Transform} (EDT) --- una mappa
delle distanze dalla costa, usata sia per i vincoli di spawn sia per la penalit\`a di
prossimit\`a alla terra. In pratica, caricamento e interrogazione del campo si riducono a:

\begin{lstlisting}
field = NetCDFLoader(data_dir).load("CL02_V2_SRC042_Conc_10mGrid.nc")
field.set_time(t_idx)                 # seleziona il frame temporale corrente
c = field.get_concentration(x, y)     # concentrazione interpolata nel punto (x, y)
\end{lstlisting}

\paragraph{Campi di vento e corrente.}
In modo analogo, \texttt{WindData} e \texttt{CurrentData} incapsulano rispettivamente la
serie temporale di vento e il campo vettoriale di corrente, esponendo le componenti
$(u, v)$ in funzione del tempo (e della posizione, per la corrente). La sincronizzazione
temporale con il campo di concentrazione \`e mantenuta in minuti simulati.

\paragraph{Gestione centralizzata.}
Il componente \texttt{DataManager} coordina l'accesso ai dati: alla creazione
\emph{scopre automaticamente} le sorgenti disponibili dai file presenti, e su richiesta
fornisce all'ambiente il campo di concentrazione di una data coppia (sorgente, scenario)
insieme ai relativi vento e corrente. Per evitare di diventare un collo di bottiglia di
I/O durante l'addestramento --- che richiede milioni di passi e quindi continue
riletture --- il \texttt{DataManager} mantiene in memoria una \emph{cache} dei quattro
scenari di vento e corrente, eliminando gli accessi al disco dopo il primo caricamento.
Una singola chiamata restituisce il campo di concentrazione e le forzanti coerenti:

\begin{lstlisting}
# Campo di una coppia (sorgente, scenario) + forzanti della STESSA versione
field, run_id = data_manager.get_random_field_for_source_version("SRC042", "V2")
wind    = data_manager.get_wind_data_for_run(run_id, wind_mapping)       # da cache
current = data_manager.get_current_data_for_run(run_id, current_mapping) # da cache
\end{lstlisting}

\subsection{Sincronizzazione dei dati}
\label{sec:sync}

Affinch\'e l'agente percepisca un ambiente fisicamente coerente, i tre flussi di dati
--- concentrazione, vento e corrente --- devono essere allineati. Il \texttt{DataManager}
garantisce due livelli di sincronizzazione.

\paragraph{Sincronizzazione di versione.}
La versione dello scenario di vento (V0--V3) \`e codificata nel nome di ogni file (es.\
\texttt{CL02\_V2\_SRC042\_Conc\_10mGrid.nc}). Al momento di fornire un campo, il
\texttt{DataManager} estrae la versione dal nome del file di concentrazione e seleziona i
corrispondenti campi di vento e corrente della \emph{stessa} versione. In questo modo \`e
impossibile un accoppiamento errato tra concentrazione e forzanti: il plume \`e sempre
associato al vento e alla corrente che lo hanno effettivamente generato.

\paragraph{Sincronizzazione temporale.}
I tre flussi hanno cadenze di campionamento native diverse: il campo di concentrazione
MIKE21 \`e prodotto con un \emph{frame} ogni $\Delta t_c = 2\,\text{min}$, mentre vento e
corrente seguono le rispettive serie, ciascuna con il proprio passo ($\Delta t_w$ per il
vento, $\Delta t_u$ per la corrente). Se questi flussi non fossero allineati, l'agente
percepirebbe un plume che avanza a una velocit\`a incoerente con il vento e la corrente che
lo trasportano. L'ambiente li riconduce perci\`o a un unico \textbf{asse temporale
assoluto}, espresso in \emph{minuti simulati}.

\medskip\noindent\emph{Istante assoluto.} Ad ogni passo, dallo step corrente si ricava
l'istante assoluto dell'episodio come somma dell'offset del chunk di partenza e del tempo
trascorso:
\begin{equation}
  t_{\text{abs}}(k) = \underbrace{f_0 \cdot \Delta t_c}_{\text{offset del chunk}}
  \;+\; \underbrace{k \cdot \Delta t / 60}_{\text{tempo trascorso}}
  \qquad [\text{min}],
  \label{eq:tabs}
\end{equation}
dove $k$ \`e il numero di step compiuti, $\Delta t = 10\,\text{s}$ il passo dell'agente e
$f_0$ il \emph{frame} di spawn del chunk (Sezione~\ref{sec:split}). \`E essenziale che
l'offset del chunk compaia in $t_{\text{abs}}$: cos\`i, quando l'episodio parte da Q3/4, non
solo il plume ma anche il vento e la corrente sono presi al punto \emph{corrispondente} del
loro profilo temporale --- l'agente incontra lo stato idrodinamico effettivamente
contemporaneo a quello stadio di dispersione, non quello dell'inizio simulazione.

\medskip\noindent\emph{Query per flusso e interpolazione.} L'istante $t_{\text{abs}}$ viene
poi convertito nell'indice, in generale \emph{frazionario}, proprio di ciascun flusso; il
valore restituito \`e l'\textbf{interpolazione lineare} tra i due campioni adiacenti:
\begin{itemize}[leftmargin=1.4em, itemsep=2pt]
  \item \textbf{concentrazione}: l'indice di frame \`e
        $f = f_0 + (t_{\text{abs}} - f_0\,\Delta t_c)/\Delta t_c$; detta $\alpha$ la sua parte
        frazionaria, il campo restituito \`e la miscela
        $(1-\alpha)\,C_{t_0} + \alpha\,C_{t_1}$ tra i frame adiacenti $t_0=\lfloor f\rfloor$ e
        $t_1=t_0+1$ --- interpolazione temporale che si somma a quella bilineare nello spazio
        (Sezione~\ref{sec:dataloader});
  \item \textbf{vento e corrente}: l'indice \`e $t_{\text{abs}}/\Delta t_{w,u}$, con analoga
        miscela lineare tra i campioni delle rispettive serie.
\end{itemize}

\medskip\noindent\emph{Cadenza.} Poich\'e
$\Delta t_c / \Delta t = 2\,\text{min} / 10\,\text{s} = 12$, il campo di concentrazione
avanza di un frame \textbf{ogni 12 step} dell'agente; nei passi intermedi l'agente non vede
un fotogramma ``congelato'' ma una versione \emph{interpolata} del campo. Ad esempio, dopo
$11$ minuti dallo spawn ($k = 66$) l'indice di frame vale $f_0 + 5{,}5$: il campo restituito
\`e la media al $50\%$ fra il frame $f_0{+}5$ e il frame $f_0{+}6$
(Figura~\ref{fig:timeline}).

\begin{figure}[H]
\centering
\begin{tikzpicture}[font=\footnotesize, x=2.7cm]
  \draw[-{Latex[length=2mm]}] (-0.15,0) -- (4.55,0) node[right]{$t$ [min]};
  % frame di concentrazione (tick alti)
  \foreach \x/\lab in {0/$f_0$, 2/$f_0{+}5$, 4/$f_0{+}6$}{
    \draw[thick] (\x,-0.12) -- (\x,0.36);
    \node[above, font=\scriptsize] at (\x,0.34) {\lab};
  }
  \node[below, font=\scriptsize] at (2,-0.12){$10$}; \node[below, font=\scriptsize] at (4,-0.12){$12$};
  % step dell'agente (tick corti tra frame f0+5 e f0+6)
  \foreach \i in {1,...,11}{ \draw[gray] ({2+\i/6},-0.06) -- ({2+\i/6},0.06); }
  % query a t=11 min (indice f0+5.5)
  \draw[densely dashed, red!70!black] (3,-0.16) -- (3,0.62);
  \node[red!70!black, above, font=\scriptsize, align=center] at (3,0.6){$k{=}66$\\($11$\,min)};
  \draw[{Latex[length=1.5mm]}-{Latex[length=1.5mm]}, blue!55!black] (2,0.18) -- (3,0.18)
       node[midway, above, font=\scriptsize]{$\alpha$};
  \draw[{Latex[length=1.5mm]}-{Latex[length=1.5mm]}, blue!55!black] (3,0.18) -- (4,0.18)
       node[midway, above, font=\scriptsize]{$1-\alpha$};
  \node[align=center, font=\scriptsize, gray!50!black] at (3,-0.62)
       {$12$ step dell'agente ($\Delta t{=}10$\,s) per ogni frame ($\Delta t_c{=}2$\,min)};
\end{tikzpicture}
\caption{Sincronizzazione temporale della concentrazione. Tra due frame consecutivi
         ($f_0{+}5$ a $10\,\text{min}$, $f_0{+}6$ a $12\,\text{min}$) l'agente compie 12 step;
         allo step $k{=}66$ ($11\,\text{min}$) il campo \`e la miscela lineare
         $(1-\alpha)C_{f_0+5} + \alpha C_{f_0+6}$ con $\alpha = 0{,}5$. Vento e corrente sono
         interpolati in modo analogo sulle rispettive serie.}
\label{fig:timeline}
\end{figure}

\medskip\noindent\emph{Gestione dei bordi.} Un episodio dura al pi\`u $1080$ step
($= 180\,\text{min}$); partendo da un chunk avanzato, $t_{\text{abs}}$ pu\`o oltrepassare la
fine di una serie. In tal caso l'indice viene \textbf{saturato} all'ultimo campione
disponibile (con un avviso una tantum per il vento), cos\`i che l'episodio resti sempre ben
definito anche quando ``sopravvive'' ai dati. In codice, l'aggiornamento dei tre flussi ad
ogni step si riduce a:

\begin{lstlisting}
dt_c  = 2.0                                       # min per frame di concentrazione
t_abs = f0 * dt_c + step * dt / 60.0              # istante assoluto in minuti simulati
field.set_time(f0 + (step * dt / 60.0) / dt_c)    # frame float -> interp. temporale lineare
wind.set_time_from_minutes(t_abs)                 # indice = t_abs / dt_wind  (clamp agli estremi)
current.set_time_from_minutes(t_abs)              # indice = t_abs / dt_curr  (clamp agli estremi)
\end{lstlisting}
